\begin{ExerciseBlock}
\ExerciseHeading{01}{Fotoeffekt bei Zink}{quantum_optics_photoelectric_effect}
Beim äußeren Fotoeffekt können Photonen Elektronen aus einem Material lösen, wenn ihre Energie mindestens die Austrittsarbeit $W_\mathrm A$ erreicht. Für das hier betrachtete Zink nehmen wir $W_\mathrm A=\qty{4.7}{\electronvolt}$ an.
\begin{enumerate}
\item Geben Sie $W_\mathrm A$ in Joule an.
\item Berechnen Sie die Grenzwellenlänge $\lambda_\mathrm{grenz}$.
\item Für $\lambda=\qty{100}{\nano\metre}$: Bestimmen Sie die Photonenenergie in Elektronenvolt und die maximale kinetische Energie der ausgelösten Elektronen.
\item Zeigen Sie, warum bei $\lambda=\qty{700}{\nano\metre}$ kein Einphotonen-Fotoeffekt auftritt.
\end{enumerate}
\end{ExerciseBlock}
\begin{Solution}
\par\medskip\noindent\textbf{Musterlösung}\par\nopagebreak[4]
\begin{enumerate}
\item Mit $\qty{1}{\electronvolt}=\qty{1.602176634e-19}{\joule}$ ergibt sich $W_\mathrm A\approx\qty{7.53e-19}{\joule}$.
\item An der Schwelle gilt $hc/\lambda_\mathrm{grenz}=W_\mathrm A$. Daher ist $\lambda_\mathrm{grenz}=hc/W_\mathrm A\approx\qty{264}{\nano\metre}$.
\item Aus $E_\mathrm{ph}=hc/\lambda$ folgt $E_\mathrm{ph}\approx\qty{12.4}{\electronvolt}$. Die schnellsten Elektronen haben $E_{\mathrm{kin,max}}=E_\mathrm{ph}-W_\mathrm A\approx\qty{7.7}{\electronvolt}$.
\item Bei \qty{700}{\nano\metre} beträgt $E_\mathrm{ph}=hc/\lambda\approx\qty{1.77}{\electronvolt}$. Das ist weniger als $W_\mathrm A$.
\end{enumerate}
\end{Solution}
